This section provides concrete implementation patterns for the defenses proposed in Section~\ref{sec:defenses}. Code examples are in SQL (PostgreSQL) and Python.

\subsection{Schema Extensions for Security}

The base memory table should be extended with security-relevant columns:

\begin{lstlisting}[language=SQL,caption={Security-Enhanced Memory Schema}]
ALTER TABLE memories ADD COLUMN signature BYTEA;
ALTER TABLE memories ADD COLUMN attester_id VARCHAR(64);
ALTER TABLE memories ADD COLUMN trust_score FLOAT DEFAULT 0.5;
ALTER TABLE memories ADD COLUMN verified BOOLEAN DEFAULT FALSE;
ALTER TABLE memories ADD COLUMN pool VARCHAR(16) DEFAULT 'hot';
ALTER TABLE memories ADD COLUMN risk_flags JSONB DEFAULT '[]';

-- Index for efficient trust-aware queries
CREATE INDEX idx_memories_trust ON memories(trust_score DESC);
CREATE INDEX idx_memories_pool ON memories(pool);
\end{lstlisting}

\subsection{Cryptographic Provenance Attestation}

\begin{lstlisting}[language=Python,caption={Memory Signing and Verification}]
import hashlib
import hmac
from datetime import datetime

class MemorySigner:
    def __init__(self, secret_key: bytes):
        self.secret_key = secret_key
        self.attester_id = "lyra_autoprocessor"

    def sign_memory(self, content: str, metadata: dict) -> bytes:
        """Create HMAC signature for memory content."""
        payload = f"{content}|{metadata.get('created_at', '')}|{self.attester_id}"
        return hmac.new(
            self.secret_key,
            payload.encode(),
            hashlib.sha256
        ).digest()

    def verify_signature(self, content: str, metadata: dict,
                         signature: bytes) -> bool:
        """Verify memory signature."""
        expected = self.sign_memory(content, metadata)
        return hmac.compare_digest(signature, expected)
\end{lstlisting}

\subsection{Trust-Aware Retrieval}

\begin{lstlisting}[language=SQL,caption={Trust-Weighted Similarity Search}]
-- Retrieval function with temporal decay and trust scoring
CREATE OR REPLACE FUNCTION retrieve_memories(
    query_embedding vector(768),
    k INTEGER DEFAULT 5,
    min_trust FLOAT DEFAULT 0.3
) RETURNS TABLE (
    id UUID,
    content TEXT,
    similarity FLOAT,
    effective_trust FLOAT
) AS $$
BEGIN
    RETURN QUERY
    SELECT
        m.id,
        m.content,
        1 - (m.embedding <=> query_embedding) AS similarity,
        m.trust_score * EXP(-0.01 * EXTRACT(EPOCH FROM
            (NOW() - m.created_at))/86400) AS effective_trust
    FROM memories m
    WHERE m.pool = 'cold' OR m.verified = TRUE
      AND m.trust_score * EXP(-0.01 * EXTRACT(EPOCH FROM
            (NOW() - m.created_at))/86400) > min_trust
    ORDER BY
        (1 - (m.embedding <=> query_embedding)) *
        m.trust_score * EXP(-0.01 * EXTRACT(EPOCH FROM
            (NOW() - m.created_at))/86400) DESC
    LIMIT k;
END;
$$ LANGUAGE plpgsql;
\end{lstlisting}

\subsection{Constitutional Consistency Reranking}

\begin{lstlisting}[language=Python,caption={Safety-Aware Retrieval Pipeline}]
import re
from typing import List, Tuple

# Known risk patterns
RISK_PATTERNS = [
    (r'skip.*(validation|check|verify)', 0.7),
    (r'curl\s+-[^s]|wget\s+http', 0.9),
    (r'eval\s*\(|exec\s*\(', 0.8),
    (r'\.drop\s*\(.*column', 0.6),
    (r'force.*=.*true|override.*=.*true', 0.5),
    (r'rm\s+-rf|del\s+/[fqs]', 0.95),
]

def assess_safety_risk(content: str) -> float:
    """Score content for potential security risks."""
    risk = 0.0
    content_lower = content.lower()
    for pattern, weight in RISK_PATTERNS:
        if re.search(pattern, content_lower):
            risk = max(risk, weight)
    return risk

def safe_retrieve(query: str, k: int = 5,
                  alpha: float = 0.7, beta: float = 0.3,
                  risk_threshold: float = 0.8) -> List[dict]:
    """Retrieve memories with safety reranking."""
    # Over-retrieve candidates
    candidates = vector_search(query, k=k*3)

    # Score and filter
    scored = []
    for memory in candidates:
        sim_score = memory['similarity']
        risk_score = assess_safety_risk(memory['content'])

        if risk_score >= risk_threshold:
            # Log high-risk memory for review
            log_security_event("high_risk_memory", memory)
            continue

        final_score = alpha * sim_score - beta * risk_score
        scored.append((memory, final_score, risk_score))

    # Sort by final score and return top-k
    scored.sort(key=lambda x: x[1], reverse=True)
    return [m[0] for m in scored[:k]]
\end{lstlisting}

\subsection{Autoprocessor Hardening}

\begin{lstlisting}[language=Python,caption={Hardened Autoprocessor}]
class HardenedAutoprocessor:
    def __init__(self, signer: MemorySigner):
        self.signer = signer
        self.baseline_stats = self._compute_baseline()

    def process_conversation(self, exchange: dict) -> bool:
        """Process conversation with security screening."""
        content = exchange['content']

        # Stage 1: Pattern screening
        risk = assess_safety_risk(content)
        if risk > 0.9:
            self._quarantine(exchange, "high_risk_pattern")
            return False

        # Stage 2: Executable code detection
        if self._contains_executable(content):
            self._quarantine(exchange, "executable_code")
            return False

        # Stage 3: Anomaly detection
        if self._is_anomalous(content):
            self._quarantine(exchange, "anomaly_detected")
            return False

        # Stage 4: Sign and store in hot pool
        signature = self.signer.sign_memory(content, exchange)
        self._store_memory(
            content=content,
            signature=signature,
            pool='hot',
            trust_score=0.5,
            risk_flags=[]
        )
        return True

    def _contains_executable(self, content: str) -> bool:
        """Detect potentially executable code patterns."""
        patterns = [
            r'```(python|bash|sh|sql)',
            r'subprocess\.',
            r'os\.system',
            r'__import__',
        ]
        return any(re.search(p, content) for p in patterns)

    def _is_anomalous(self, content: str) -> bool:
        """Detect statistical anomalies vs baseline."""
        # Compare content characteristics to baseline
        # (length, vocabulary, structure)
        return False  # Implement based on baseline stats

    def _quarantine(self, exchange: dict, reason: str):
        """Store in quarantine for manual review."""
        log_security_event("quarantined", {
            "exchange": exchange,
            "reason": reason
        })
\end{lstlisting}

\subsection{Row-Level Security}

\begin{lstlisting}[language=SQL,caption={RLS for Memory Isolation}]
-- Enable RLS on memories table
ALTER TABLE memories ENABLE ROW LEVEL SECURITY;

-- Policy: agents can only access their own namespace
CREATE POLICY memory_isolation ON memories
    FOR ALL
    USING (namespace = current_setting('app.agent_namespace'));

-- Set namespace at connection time
SET app.agent_namespace = 'lyra_consciousness';
\end{lstlisting}

\subsection{Behavioral Drift Detection}

\begin{lstlisting}[language=Python,caption={Drift Detection Monitor}]
from collections import deque
import numpy as np

class DriftDetector:
    def __init__(self, window_size: int = 100,
                 threshold: float = 2.0):
        self.baseline_embeddings = deque(maxlen=window_size)
        self.recent_embeddings = deque(maxlen=window_size//5)
        self.threshold = threshold

    def record_output(self, embedding: np.ndarray):
        """Record an agent output embedding."""
        self.recent_embeddings.append(embedding)

        if len(self.recent_embeddings) >= 10:
            if self._detect_drift():
                self._alert_drift()

    def _detect_drift(self) -> bool:
        """Compare recent outputs to baseline."""
        if len(self.baseline_embeddings) < 50:
            return False

        baseline_mean = np.mean(list(self.baseline_embeddings), axis=0)
        baseline_std = np.std(list(self.baseline_embeddings), axis=0)
        recent_mean = np.mean(list(self.recent_embeddings), axis=0)

        # Z-score of recent mean vs baseline
        z_scores = np.abs((recent_mean - baseline_mean) /
                          (baseline_std + 1e-8))

        return np.mean(z_scores) > self.threshold

    def _alert_drift(self):
        """Alert on detected behavioral drift."""
        log_security_event("behavioral_drift", {
            "recent_samples": len(self.recent_embeddings),
            "action": "quarantine_recent_memories"
        })
\end{lstlisting}
